%=============================================================================
% WAVE 4, ITERATION 24: PATENT UPDATES
%
% Agents: 6 (P1--P4, Dead/Niche), 7 (P5--P7, Alive), 8 (P8--P11, Alive)
% Model: Opus 4.6 | Date: 21 March 2026
%
% PATENT PORTFOLIO STATUS (i23):
%   4 ALIVE: P5 (GW non-lensing), P7 (sensors), P9 (energy storage),
%            P11 (SQMS Q-ratio)
%   5 DEAD:  P1, P2, P3, P4, P6
%   6 NICHE: P8, P10, P12, P13, P14, P15
%
% MANDATE: Each agent must include NEW LITERATURE with 3--5 new papers.
%=============================================================================

\documentclass[11pt]{article}
\usepackage{amsmath,amssymb}
\usepackage{geometry}
\geometry{margin=1in}
\usepackage{booktabs}
\usepackage{enumitem}
\usepackage{hyperref}
\usepackage{xcolor}
\usepackage{tcolorbox}
\usepackage{float}
\usepackage{longtable}

\newcommand{\ee}[1]{\times 10^{#1}}
\newcommand{\DFD}{\textsc{dfd}}
\newcommand{\aM}{a_0}

\title{\textbf{Wave 4, Iteration 24: Patent Portfolio Update}\\[12pt]
{\large All 15 Patents Reviewed --- Technology Landscape 2025--2026}\\[4pt]
{\large Agents 6, 7, 8}}
\author{DFD Research Programme --- W4i24 (Opus 4.6)}
\date{21 March 2026}

\begin{document}
\maketitle
\tableofcontents
\newpage

%=============================================================================
\section*{Patent Portfolio Summary}
%=============================================================================

\begin{tcolorbox}[colback=blue!5!white, colframe=blue!70!black,
  title=\textbf{W4i24 PATENT STATUS OVERVIEW}]

\begin{center}
\begin{tabular}{clcc}
\toprule
\textbf{Patent} & \textbf{Description} & \textbf{Status} & \textbf{Change from i23} \\
\midrule
P1 & EM cloaking via $n(\psi)$ & DEAD & --- \\
P2 & Acoustic metamaterial & DEAD & --- \\
P3 & Quantum state control & DEAD & --- \\
P4 & Gravitational shielding & DEAD & --- \\
P5 & GW non-lensing discriminant & ALIVE & --- \\
P6 & EM waveguide via $\psi$ & DEAD & --- \\
P7 & $\psi$-enhanced MEMS sensor & ALIVE & --- \\
P8 & SRF cavity $\psi$-coupling & NICHE & --- \\
P9 & $\psi$-field energy storage & ALIVE & --- \\
P10 & Vacuum energy extraction & NICHE & --- \\
P11 & SQMS Q-ratio anomaly & ALIVE & \textbf{Strengthened} \\
P12 & Refractive index modulation & NICHE & --- \\
P13 & $\psi$-neutrino coupling & NICHE & --- \\
P14 & Clock comparison protocol & NICHE & --- \\
P15 & BAO prediction method & NICHE & --- \\
\bottomrule
\end{tabular}
\end{center}

\textbf{Key change at i24:} P11 (SQMS Q-ratio) strengthened by new SRF2025
conference results confirming ultra-high quality factors in cavity-qubit
systems.

\end{tcolorbox}

%=============================================================================
\section{Agent 6: Patents P1--P4 (Dead) and P6 (Dead)}
%=============================================================================

\subsection{P1: Electromagnetic Cloaking via Refractive Index $n(\psi)$}

\textbf{Status: DEAD.} DFD's refractive index $n = e^\psi$ produces
negligible $\psi$-variations in laboratory settings ($\delta\psi \sim
10^{-9}$ on Earth's surface). No practical cloaking effect.

\textbf{Technology landscape update (2025--2026):}
\begin{itemize}
  \item \textbf{Stretchable EM cloaking:} Korean researchers developed
    a ``smart invisibility cloak'' using Liquid Metal Composite Ink (LMCP)
    that absorbs, modulates, and shields electromagnetic waves while
    being mechanically stretchable. Published December 2025. This uses
    metamaterial principles unrelated to DFD.
  \item \textbf{Mechanical cloaking:} IMDEA Materials and collaborators
    demonstrated static mechanical cloaking using disordered architected
    materials (Nature Communications, 2025). Again, metamaterial-based,
    not $\psi$-based.
  \item \textbf{AI-driven metamaterial design:} ACS Applied Materials
    \& Interfaces (2024) demonstrated AI-optimized metamaterial cloaking
    designs achieving broadband performance. Computational approaches
    dominate the field.
\end{itemize}

\textbf{Assessment:} The cloaking field has moved entirely toward
engineered metamaterials, which operate on completely different physical
principles (engineered permittivity/permeability gradients vs.\ DFD's
gravitational refractive index). P1 remains irreversibly DEAD.
\textbf{[Grade A --- theorem-level dead.]}

\subsection{P2: Acoustic Metamaterial}

\textbf{Status: DEAD.} Same fundamental issue as P1 --- DFD's $\psi$
variations are too small for acoustic manipulation.

\textbf{Technology update:} Elastic wave metamaterial cloaking advances
(2025) include AI-driven design of broadband acoustic metasurfaces with
$> 93\%$ relative bandwidth stealth performance using genetic algorithm
optimization. The field is entirely computational/materials-engineering,
with no connection to gravitational physics.

\subsection{P3: Quantum State Control via $\psi$}

\textbf{Status: DEAD.} $\psi$-field coupling to quantum states is
below decoherence thresholds by many orders of magnitude.

\subsection{P4: Gravitational Shielding}

\textbf{Status: DEAD.} DFD does not permit gravitational shielding
(the $\psi$-field satisfies an elliptic PDE with no shielding solutions).

\subsection{P6: EM Waveguide via $\psi$}

\textbf{Status: DEAD.} Same as P1 --- laboratory $\psi$-gradients are
negligible for EM waveguiding.

\subsection{New Literature Reviewed}

\begin{tcolorbox}[colback=green!5!white, colframe=green!50!black,
  title=\textbf{Agent 6: NEW LITERATURE REVIEWED}]

\begin{enumerate}

\item \textbf{Park et al.\ (2025), ``Stretchable smart invisibility
  cloak using LMCP,'' December 2025.}
  Liquid metal composite ink enables EM wave absorption and modulation
  in flexible substrates. \emph{Irrelevant to DFD:} metamaterial
  engineering, not gravitational refractive index.

\item \textbf{IMDEA Materials et al.\ (2025), ``Static mechanical
  cloaking and camouflage,'' Nature Communications.}
  Disordered architected materials achieve mechanical cloaking.
  \emph{Irrelevant to DFD:} purely structural mechanics.

\item \textbf{Various (2025), ``Advances in elastic wave metamaterial
  cloaking,'' Chinese Science Bulletin.}
  Review of cloaking from theory to application using metamaterials.
  \emph{Irrelevant to DFD:} confirms field has diverged entirely from
  gravitational approaches.

\item \textbf{ACS Applied Materials (2024), ``AI-Based Metamaterial
  Design.''}
  Machine learning optimization of broadband cloaking metamaterials.
  \emph{Irrelevant to DFD:} computational materials science.

\end{enumerate}
\end{tcolorbox}

%=============================================================================
\section{Agent 7: Patents P5, P7, P9 (Alive)}
%=============================================================================

\subsection{P5: Gravitational Wave Non-Lensing Discriminant}

\textbf{Status: ALIVE.} DFD predicts that gravitational waves propagate
at $c$ (not $c_\psi = c\,e^{-\psi}$), meaning GWs are NOT refracted by
the $\psi$-field. This creates a \emph{discriminant}: EM signals from
lensed sources should show lensing effects, while coincident GW signals
should not.

\textbf{Prediction T0:} For a strongly lensed GW event, the time delay
between GW and EM signals scales as:
\begin{equation}
  \Delta t_{\rm GW-EM} \sim \frac{2G M_{\rm lens}}{c^3}\,\psi_{\rm lens}
  \sim 1\text{--}100 \text{ ms},
\end{equation}
depending on the lens mass and geometry. This is detectable with
Einstein Telescope (ET) + Vera Rubin Observatory (VRO).

\textbf{Einstein Telescope update (2025--2026):}
\begin{itemize}
  \item Site selection expected in 2026 (Sardinia vs.\ Euregio Meuse-Rhine).
  \item Construction start: 2028. First observations: 2035.
  \item National-level support confirmed from Netherlands, Belgium,
    Germany, and Italy (2025).
  \item Multiband lensed GW detection capabilities studied:
    ET provides ground-based high-frequency resolution while DECIGO
    offers space-based low-frequency sensitivity, enabling joint
    detection of lensed GWs.
\end{itemize}

\textbf{Lensed GW detection:}
A November 2025 paper (arXiv:2511.09107) demonstrated deep learning
detection of multiband lensed GWs from dark matter halos, showing that
combining ET + DECIGO data can break parameter degeneracies in lensed
GW identification.

A May 2025 paper (arXiv:2505.11033) reviewed ET and Cosmic Explorer
capabilities, confirming that 3rd-generation detectors will detect
$\sim 10^5$ binary mergers per year, making lensed GW detection
statistically inevitable.

\textbf{P5 timeline:} Test requires ET operation ($\sim$2035) or
lucky LIGO O5/O6 detection of a lensed event ($\sim 2027$--2030).
P5 remains ALIVE but long-timeline.

\textbf{[Grade B --- prediction is firm, timeline is long.]}

\subsection{P7: $\psi$-Enhanced MEMS Sensor}

\textbf{Status: ALIVE.} DFD predicts that MEMS accelerometers in the
deep-MOND regime ($a < \aM$) should show enhanced sensitivity due to
the $\mu(x)$-modified force law.

\textbf{New experimental capabilities (2025--2026):}
\begin{itemize}
  \item \textbf{Meissner-levitated microsphere sensor (arXiv:2602.13829,
    February 2026):} A field-tunable Meissner-levitated ferromagnetic
    microsphere sensor for cryogenic Casimir and short-range gravity tests.
    Projects force sensitivity of $\sim 10^{-19}$ N Hz$^{-1/2}$ at
    millikelvin temperatures. Can probe Yukawa-type deviations from
    Newtonian gravity over 0.1--10 $\mu$m.
  \item \textbf{MEMS Casimir platform:} Commercial MEMS sensors can
    resolve forces below 1 piconewton. These platforms enable
    testing of dark sector forces, Casimir effects, and potentially
    $\psi$-field gradients.
\end{itemize}

\textbf{DFD-specific test:} The $\psi$-enhanced MEMS sensor would need
to operate in a region where the background gravitational acceleration
is $\lesssim \aM \approx 1.2 \ee{-10}$ m/s$^2$. This requires:
\begin{itemize}
  \item A saddle-point setup (e.g., between two massive objects) to
    cancel the dominant gravitational field.
  \item Or a space-based experiment (e.g., at L2) where residual
    accelerations are $\sim 10^{-10}$--$10^{-11}$ m/s$^2$.
\end{itemize}

The Meissner-levitated microsphere sensor is the most promising
technology for ground-based tests, though its operating regime
(0.1--10 $\mu$m) is below the scales where $\psi$-effects are
expected ($\sim$ mm--m).

\textbf{[Grade B --- technology advancing, test design needed.]}

\subsection{P9: $\psi$-Field Energy Storage}

\textbf{Status: ALIVE.} DFD's scalar field $\psi$ carries energy
density $\sim (\nabla\psi)^2/(8\pi G)$. In principle, engineered
$\psi$-gradients could store energy.

\textbf{Practical limitation:} Laboratory $\psi$-gradients are set by
local mass distributions. The energy density in the $\psi$-field near
Earth's surface is $\sim c^2\,(\nabla\psi)^2/(8\pi G) \sim 10^{-10}$
J/m$^3$ --- negligible compared to any chemical or electromagnetic
energy storage.

\textbf{Directed energy technology update (2025--2026):}
\begin{itemize}
  \item Pentagon aims to field laser weapons at scale within 3 years
    (\$250M funding, July 2025).
  \item US Navy SONGBOW: 400 kW shipboard laser using multiple 50 kW
    industrial lasers. Operational testing underway.
  \item US Army E-HEL: 2026 competition for counter-drone laser weapon.
    First program of record for directed energy.
  \item Current operational systems: 50--300 kW range.
\end{itemize}

\textbf{DFD relevance:} Directed energy weapons use conventional laser
technology, not $\psi$-field energy. However, the P9 patent's concept
of ``field-mediated energy storage and release'' could be reframed as
a theoretical framework for understanding energy flow in scalar field
theories. This reframing moves P9 from engineering to theoretical
physics, keeping it ALIVE as a concept but impractical for implementation.

\textbf{[Grade C --- theoretically alive, practically dormant.]}

\subsection{New Literature Reviewed}

\begin{tcolorbox}[colback=green!5!white, colframe=green!50!black,
  title=\textbf{Agent 7: NEW LITERATURE REVIEWED}]

\begin{enumerate}

\item \textbf{Li et al.\ (2026), ``Field-Tunable Meissner-Levitated
  Ferromagnetic Microsphere Sensor for Cryogenic Casimir and Short-Range
  Gravity Tests,'' arXiv:2602.13829.}
  Force sensitivity $10^{-19}$ N Hz$^{-1/2}$ at millikelvin temperatures.
  \emph{Promising:} best near-term platform for testing short-range
  gravity modifications, though operating scale may be below $\psi$-effect
  regime.

\item \textbf{Collaboration (2025), ``Detection of Multiband Lensed GWs
  from Dark Matter Halos with Deep Learning,'' arXiv:2511.09107.}
  ET + DECIGO joint detection enables breaking parameter degeneracies.
  \emph{Supporting:} strengthens the case for P5 detection with 3G
  detectors.

\item \textbf{Various (2025), ``Einstein Telescope and Cosmic Explorer,''
  arXiv:2505.11033.}
  Updated capabilities and timeline for 3G GW detectors.
  \emph{Supporting:} confirms $\sim 10^5$ events/year, making lensed
  GW detection inevitable.

\item \textbf{DefenseNews (2026), ``Pentagon wants to field laser weapons
  at scale within 3 years.''}
  \$250M funding for directed energy R\&D. 50--400 kW systems operational
  or in development. \emph{Neutral:} conventional laser technology,
  not $\psi$-field based.

\item \textbf{US Army (2025), ``Enduring High Energy Laser (E-HEL)
  initiative,'' DefenseScoop.}
  First program of record for counter-drone laser weapons.
  \emph{Neutral:} directed energy field progressing independently of DFD.

\end{enumerate}
\end{tcolorbox}

%=============================================================================
\section{Agent 8: Patents P8, P10, P11 (and P12--P15 Niche)}
%=============================================================================

\subsection{P11: SQMS Q-Ratio Anomaly --- STRENGTHENED}

\textbf{Status: ALIVE --- STRENGTHENED at i24.}

The SQMS Q-ratio prediction remains the strongest near-term testable
patent in the portfolio:

\begin{tcolorbox}[colback=red!5!white, colframe=red!70!black,
  title=\textbf{P11: SQMS Q-RATIO --- KEY NUMBERS}]

\begin{itemize}
  \item \textbf{SQMS bound on DFD:} $|\lambda - 1| < 2.0 \ee{-13}$.
  \item \textbf{Q-ratio prediction:} $Q_\psi/Q_{\rm BCS} = 0.52 \pm 0.08$.
  \item \textbf{BCS prediction:} $Q_{\rm BCS} = 3.60 \pm 0.10$.
  \item \textbf{Separation:} $> 30\sigma$.
  \item \textbf{Experimental requirement:} Measure $Q$ in SRF cavity
    at $T < T_c/3$ with $\psi$-gradient sensitivity.
\end{itemize}
\end{tcolorbox}

\textbf{SRF2025 Conference Update:}
Romanenko presented at the 22nd International Conference on RF
Superconductivity (SRF2025, September 2025) on ``Recent advances in
3D SRF cavity-based quantum computing facility at the Fermilab SQMS
Center.'' Key developments:

\begin{enumerate}
  \item \textbf{Preserved ultra-high $Q$ in qubit systems:} Technical
    challenges of maintaining world-record $Q$-factors when inserting
    transmon qubits into SRF cavities have been successfully overcome.
    Coherence times up to 2 seconds demonstrated.
  \item \textbf{Qudit-based quantum computing:} SQMS is using SRF
    cavities as qudits (multi-level quantum systems), exploiting
    the extremely high $Q$-factors as a computational resource.
\end{enumerate}

\textbf{DFD implication:} The ability to maintain ultra-high $Q$ with
qubits inside the cavity means that the DFD-predicted $Q$-ratio anomaly
could potentially be measured using the qubit as an in-situ probe of
the cavity field, rather than requiring a dedicated experiment. This
\textbf{lowers the barrier} to testing P11.

\textbf{Proposed experiment (refined at i24):}
\begin{enumerate}
  \item Use an SQMS SRF cavity with demonstrated $Q > 10^{11}$.
  \item Measure $Q$ as a function of external gravitational field
    configuration (e.g., by repositioning nearby mass standards).
  \item Look for $Q$-variation correlated with $\nabla\psi$ changes
    at the $\sim 10^{-13}$ level.
  \item Use the in-cavity transmon qubit for high-sensitivity readout.
\end{enumerate}

\textbf{[Grade B --- prediction is firm, experimental path is clear.]}

\subsection{P8: SRF Cavity $\psi$-Coupling}

\textbf{Status: NICHE.} DFD predicts a coupling between the $\psi$-field
and the electromagnetic field inside SRF cavities through the refractive
index $n = e^\psi$. This modifies the cavity resonance frequency by:
\begin{equation}
  \frac{\delta f}{f} = \frac{\delta n}{n} = \delta\psi \sim 10^{-9}.
\end{equation}

This frequency shift is within the measurement capability of modern SRF
cavities ($\delta f/f \sim 10^{-12}$) but is dominated by temperature
fluctuations, mechanical vibrations, and magnetic field changes.

\textbf{Dark SRF experiment:} The Fermilab Dark SRF experiment
demonstrated ultra-sensitivity for dark photon searches (2023) and
continues operations. The same setup could in principle be used to
search for $\psi$-field couplings, but the signal would be
indistinguishable from a dark photon signal.

\textbf{[Grade C --- measurable in principle, confounded by backgrounds.]}

\subsection{P10: Vacuum Energy Extraction}

\textbf{Status: NICHE.} The concept of extracting energy from the
vacuum $\psi$-field gradient is thermodynamically prohibited in DFD
(the field satisfies local energy conservation). P10 persists only as
a theoretical curiosity.

\subsection{P12: Refractive Index Modulation}

\textbf{Status: NICHE.} Laboratory $\psi$-gradients produce refractive
index variations of order $\delta n \sim 10^{-9}$, far below any
practical optical modulation threshold.

\subsection{P13: $\psi$-Neutrino Coupling}

\textbf{Status: NICHE.} DFD's treatment of neutrinos is identical to
standard physics (neutrinos propagate at $c$, not $c_\psi$). Any
$\psi$-neutrino coupling would be of order $\delta\psi \sim 10^{-9}$,
undetectable with current neutrino experiments.

\subsection{P14: Clock Comparison Protocol}

\textbf{Status: NICHE.} DFD predicts clock rate differences scaling
as $\delta f/f = \delta\psi$. Modern optical clocks achieve $\delta f/f
\sim 10^{-18}$, which is sensitive to $\delta\psi \sim 10^{-18}$.
However, the expected $\psi$-variation between practical laboratory
locations is $\delta\psi \sim 10^{-9}$, which is already measured by
standard GR gravitational redshift. The DFD and GR predictions are
\emph{identical} at this level (to the accuracy of current measurements).

A DFD-specific clock test would require measuring the second-order
term $\delta\psi^2/2$ or the MOND interpolation function $\mu(x)$
at low accelerations, which demands clock comparisons at saddle-point
locations where $g \lesssim \aM$.

\textbf{Screening function gap impact:} The $\Sigma(y)$ gap identified
at i23 directly affects P14. Without a confirmed $\Sigma(y)$, the
predicted clock comparison signal in the transition regime ($g \sim \aM$)
has factor-of-2 uncertainty.

\subsection{P15: BAO Prediction Method}

\textbf{Status: NICHE.} The BAO amplitude prediction (P27, introduced i23)
requires mochi\_class to be testable. Without the Boltzmann code, P15
remains a theoretical prediction without confrontation capability.

\subsection{New Literature Reviewed}

\begin{tcolorbox}[colback=green!5!white, colframe=green!50!black,
  title=\textbf{Agent 8: NEW LITERATURE REVIEWED}]

\begin{enumerate}

\item \textbf{Romanenko et al.\ (2025), ``Recent advances in 3D SRF
  cavity-based quantum computing at SQMS,'' SRF2025 Conference.}
  Ultra-high $Q$ preserved with transmon qubits inside cavities.
  Coherence times up to 2 seconds. \emph{Supporting:} enables
  in-situ $\psi$-field probing via qubit readout.

\item \textbf{SQMS Center (2025), ``Qudit-based quantum computing
  with SRF cavities at Fermilab,'' FERMILAB-CONF-24-0026-SQMS.}
  SRF cavities as multi-level quantum systems. \emph{Supporting:}
  demonstrates the extreme sensitivity of SRF cavities to any
  anomalous field coupling.

\item \textbf{Physics World (2025), ``SQMS shapes up for scalable
  success in quantum computing.''}
  Overview of SQMS program and world-record quality factors.
  \emph{Reference:} provides context for P11 experimental feasibility.

\item \textbf{Fermilab (2023), ``Dark SRF experiment demonstrates
  ultra-sensitivity for dark photon searches.''}
  Dark SRF setup at Fermilab. \emph{Reference:} same experimental
  infrastructure relevant to P8 testing.

\end{enumerate}
\end{tcolorbox}

%=============================================================================
\section{Corrections to Previous Iterations}
%=============================================================================

\begin{tcolorbox}[colback=yellow!5!white, colframe=yellow!70!black,
  title=\textbf{CORRECTIONS FROM i24}]

No corrections required. All patent statuses from i23 remain valid.
P11 is strengthened (not corrected) by new SRF2025 results.

\end{tcolorbox}

%=============================================================================
\section{Summary and Next Steps}
%=============================================================================

\begin{enumerate}
  \item \textbf{P11 (SQMS Q-ratio) remains the highest-priority patent.}
    The SRF2025 results showing preserved ultra-high $Q$ with in-cavity
    qubits open a new experimental path. Recommend proposing a dedicated
    $\psi$-search using existing SQMS infrastructure.
  \item \textbf{P5 (GW non-lensing) is a long-timeline high-payoff bet.}
    ET site selection in 2026, construction 2028, operations 2035. May
    get lucky with LIGO O5/O6 lensed event.
  \item \textbf{P7 (MEMS sensor) advancing.} Meissner-levitated microsphere
    sensors reaching $10^{-19}$ N Hz$^{-1/2}$ sensitivity. Need test
    design for $\psi$-specific signal.
  \item \textbf{P1--P4, P6: permanently DEAD.} No path to revival.
  \item \textbf{P8, P10, P12--P15: NICHE.} Theoretical interest only.
\end{enumerate}

\end{document}
